% Generated draft: requires mathematical and source review.
\documentclass[11pt]{article}
\usepackage[margin=25mm]{geometry}
\usepackage{fontspec}
\setmainfont{Latin Modern Roman}
\usepackage{amsmath,amssymb,mathtools,unicode-math}
\setmathfont{Latin Modern Math}
\usepackage{enumitem,needspace,xurl,hyperref}
\setlength{\parindent}{0pt}
\setlength{\parskip}{5pt}
\setlength{\emergencystretch}{3em}
\newcommand{\N}{\mathbb{N}}
\newcommand{\Z}{\mathbb{Z}}
\newcommand{\Q}{\mathbb{Q}}
\newcommand{\R}{\mathbb{R}}
\newcommand{\C}{\mathbb{C}}
\newcommand{\F}{\mathbb{F}}
\newcommand{\A}{\mathbb{A}}
\newcommand{\PP}{\mathbb{P}}
\newcommand{\E}{\mathbb{E}}
\newcommand{\eps}{\varepsilon}
\newcommand{\dd}{\,\mathrm{d}}
\begin{document}
\section*{UnsolvedMath ID: 3100050}
\textbf{Source problem code:} catalogue entry
\subsection{Definitions and mathematical statement}
For a real number $x\geq0$, write its binary expansion as
\[
x=\lfloor x\rfloor+\sum_{n=1}^{\infty}d_n(x)2^{-n},\qquad d_n(x)\in\{0,1\},
\]
using the expansion that is not eventually all $1$ when there are two choices. Say that $x$ has arbitrarily long zero blocks if
\[
\forall L\in\N_{\geq1}\ \exists m\in\N_{\geq1}\ \forall r\in\{0,\ldots,L-1\},\quad d_{m+r}(x)=0.
\]
The first question, attributed to Erd\H{o}s, asks whether $x=\sqrt2$ has this property. The second asks for a single algebraic number with this property. Here algebraic means a root of a nonzero polynomial in $\Z[t]$. The source does not explicitly require the second number to be real, positive, or irrational; an ordinary binary expansion requires a real-number convention, and allowing rational numbers with terminating binary expansions makes this second existence question immediate. The mathematically substantive interpretation suggested by the supplied research summary is
\[
\exists x\in(\R\setminus\Q)\text{ algebraic such that }|x|\text{ has arbitrarily long zero blocks}.
\]
The irrationality restriction in this last display is an interpretation requiring review, not a hypothesis present in the quoted question. No frequency of the blocks is requested: arbitrarily long blocks is weaker than normality in base $2$.
\subsection{Short English statement}
Does the binary expansion of $\sqrt2$ contain blocks of consecutive zeros of every finite length? The source also asks for one algebraic example, apparently intending an algebraic irrational, although it omits that restriction.
\subsection{Sources}
\begin{itemize}
\item Local UnsolvedMath registry, permanent ID 3100050, exact source code catalogue entry, attributed to Erd\H{o}s; registry status partially solved and classification PARTIAL-PROGRESS.
\item Joshua Cooper, \emph{Combinatorial Problems I Like}, October 2020 snapshot, Combinatorial Number Theory, source-order bullet 50, \url{https://people.math.sc.edu/cooper/combprob.html}. Supplied extraction: HTML, accessed July 29, 2026. The page labels the item open in that snapshot; later resolution was not checked item by item. Public access was verified in the supplied provenance, but redistribution permission was not established; both status and rights are flagged NEEDS\_REVIEW.
\item The supplied research summary reports the long-zero-block question open and connects it to normality conjectures. It names Borel, Bailey--Borwein--Crandall, and Bugeaud and others, but supplies no titles, dates, or precise citations for their results. Its assertions about strong normality and algebraic numbers require independent verification; they are not used as mathematical premises here.
\end{itemize}
\end{document}
